\subsection{Operatore posizione}
Il problema è decomporre \(\ket{\psi}\) su una base che ne dia la posizione. Allora introduco
\[
\hat{x} \ket{x} = x \ket{x}
\]
che in generale dà
\[
\braket{x|\psi} = \psi(x) \qquad \text{funzione d'onda}
\]
dove \(|\psi(x)|^2 = |\braket{x|\psi}|^2\). Dato che \(x\) è un parametro reale continuo, vorrei scrivere
\[
\ket{\psi} = \int \psi(x) \ket{x} \, dx
\]
come estensione del caso discreto. Così \(|\psi(x)|^2\) è una distribuzione di probabilità \(\longrightarrow\) una densità di probabilità del tipo
\[
P_\Delta \left( [x, x + \Delta x] \right) = \int_x^{x + \Delta x} |\psi(x)|^2 \, dx
\]
con la condizione che
\[
\int_A^B |\psi(x)|^2 \, dx = 1
\]
che esteso a tutto lo spazio diventa
\[
\int_{-\infty}^{+\infty} |\psi(x)|^2 \, dx = 1
\]
I problemi sorgono nell'esprimere la condizione di ortonormalità, perché
\[
\begin{cases}
\ket{\psi} = \int_{-\infty}^{+\infty} \ket{x} \braket{x|\psi} \, dx \\
\braket{x'|\psi} = \psi(x')
\end{cases} \quad \Rightarrow \quad \psi(x') = \int_{-\infty}^{+\infty} \braket{x'|x} \braket{x|\psi} \, dx \quad \Rightarrow
\]
\[
\Rightarrow \quad \psi(x') = \int_{-\infty}^{+\infty} \braket{x'|x} \psi(x) \, dx
\]
Cioè ho introdotto un oggetto tale che
\[
\braket{x'|x} = \delta(x',x)
\]
che ha la proprietà per definizione
\[
\int_{-\infty}^{+\infty} \psi(x) \delta(x, x') \, dx = \psi(x')
\]
Da questa discende:
\begin{enumerate}
\item invarianza per traslazione
\[
\delta(x', x) = \delta(x' - x)
\]
Poiché \(\delta(x' - x) = \delta(x' + a, x + a)\)
\[
\int_{-\infty}^{+\infty} \delta(x' + a, x + a) \psi(x) \, dx = [x'' = x + a] =
\]
\[
= \int_{-\infty}^{+\infty} \delta(x' + a, x'') \psi(x'' - a) \, dx'' = \psi(x' + a - a) = \psi(x')
\]
\item simmetria
\[
\delta(x' - x)  = \delta(x - x')
\]
Infatti
\[
\braket{\psi|x'} = \int_{-\infty}^{+\infty} \braket{\psi|x}\braket{x|x'} \, dx = \int_{-\infty}^{+\infty} \psi^*(x) \delta(x,x') \, dx = \psi^*(x')
\]
\item \(f(x) \delta(x) = f(0) \delta(x)\)
\begin{align*}
\int_{-\infty}^{+\infty} f(x) \delta(x) \psi(x) \, dx &= f(0) \psi(0) = \\
&= f(0) \int_{-\infty}^{+\infty} \psi(x) \delta(x) \, dx = \\
&= \int_{-\infty}^{+\infty} f(0) \psi(x) \delta(x) \, dx
\end{align*}
Poiché \(\psi\) è generica, \(f(0) \delta(x) = f(x) \delta(x)\).
\item \(\delta(ax) = \frac{1}{a} \delta(x)\)
\[
\int_{-\infty}^{+\infty} \delta(ax) \psi(x) \, dx = \left[ ax = t \; \Rightarrow \; dx = \frac{dt}{a} \right] = 
\]
\[
= \frac{1}{a} \int_{-\infty}^{+\infty} \delta(t) \psi \left( \frac{t}{a} \right) \, dt = \frac{1}{a} \psi(0) = \frac{1}{a} \int_{-\infty}^{+\infty} \delta(x) \psi(x) \, dx 
\]
\item normalizzazione
\[
\int_{-\infty}^{+\infty} \delta(x) \, dx = 1
\]
\[
\int_{-\infty}^{+\infty} \delta(x) \cdot 1 \, dx = 1
\]
\end{enumerate}
\textbf{Relazione di ortonormalizzazione}
\[
\ket{\psi} = \int_{-\infty}^{+\infty} \ket{x}\braket{x|\psi} \, dx
\]
Nel caso continuo \(\braket{x|x'} = \delta(x - x') \longrightarrow\) ortogonalizzazione impropria, infatti
\[
\begin{cases}
\braket{x|x} = \delta(0) \\
\int_{-\infty}^{+\infty} |\psi(x)|^2 \, dx = 1
\end{cases}
\]
\textbf{Risoluzione dell'identità}
\[
\int_{-\infty}^{+\infty} \ket{x} \bra{x} \, dx = \mathbbm{1}
\]
il cui elemento
\[
\Braket{x'|\int_{-\infty}^{+\infty} \Ket{x} \Bra{x} dx\;|x''} = \int_{-\infty}^{+\infty} \delta(x' - x) \delta(x - x'') \, dx = \delta(x' - x'') = \braket{x'|x''}
\]
\textbf{Prodotto scalare fra stati}
\[
\braket{\varphi|\psi} = \int_{+\infty}^{-\infty} \braket{\varphi|x}\braket{x|\psi} \, dx = \int_{-\infty}^{+\infty} \varphi(x) \psi(x) \, dx
\]
\[
\braket{\psi|\psi} = \int_{-\infty}^{+\infty} |\psi(x)|^2 \, dx
\]
\textbf{Operatori}
\[
A(\hat{x}) = \sum_{i = 0}^\infty a_i \, (\hat{x})^i
\]
\[
\braket{x'|A(\hat{x})|x''} = \sum_i a_i \, (\hat{x})^i \braket{x'|x''} = \delta(x' - x'') \sum_i a_i \, (\hat{x})^i = A(x'') \delta(x' - x'')
\]
Ad es.
\[
\braket{y|\cos \hat{x}|x} = \cos z \; \delta(y - x) = \cos y \; \delta(y - x)
\]
\begin{align*}
\braket{\psi|\cos \hat{x}| \psi} &= \int_{-\infty}^{+\infty} \int_{-\infty}^{+\infty} \braket{\psi|y} \braket{y|\cos \hat{x}|z} \braket{z|\psi} \, dz \, dy = \\
&= \int_{-\infty}^{+\infty} \int_{-\infty}^{+\infty} \psi^*(y) \, \cos z \, \delta(y - z) \, \psi(z) \, dy \, dz = \\
&= \int_{-\infty}^{+\infty} |\psi(z)|^2 \, \cos z \, dz
\end{align*}
\textbf{Riepilogo}
\begin{enumerate}
\item \(\braket{q'|\hat{q}|q} = q \braket{q'|q} = q \delta(q' - q)\)
\item \(\braket{q'|f(\hat{q})|q} = f(q) \braket{q|q'} = f(q) \, \delta(q - q')\)
\end{enumerate}

\subsubsection{Esempio: ulteriori proprietà della \(\delta\)}
\[
\int_{-\infty}^{+\infty} \delta(f(x)) \psi(x) \, dx
\]
\begin{align*}
\int_{-\infty}^{+\infty} \delta(f(x)) \psi(x) \, dx &= \left[ x^\prime = f(x), dx^\prime = f^\prime(x) \, dx, x = f^{-1}(x) \right] = \\
&= \int_{-\infty}^{+\infty} \frac{\delta^\prime(x)}{f^\prime(x)} \psi(f^{-1}(x)) \, dx^\prime = \\
\int_{-\infty}^{+\infty} \delta(f(x)) \psi(x) \, dx &= \sum_i \int_{a_i}^{b_i} \delta(f(x)) \psi(x) \, dx \\ 
\int_{-\infty}^{+\infty} \delta(f(x)) \psi(x) \, dx &= \frac{\psi(x_i)}{f^\prime(x)|_{x : f(x) = 0}} \psi(f^{-1})
\end{align*}

\subsubsection{Esempio: operatore parità}
\( \braket{x|\mathcal{P}|\psi} = \psi(-x) \) con \( \braket{x|\psi} = \psi(x) \Rightarrow \bra{x} \mathcal{P} = \bra{-x} \).
\begin{align*}
(\braket{\psi|\mathcal{P}|x})^* &= \left( \int_{-\infty}^{+\infty} \braket{\psi|x'} \braket{x'|\mathcal{P}|x} \, dx' \right)^* = \\
&= \left( \int_{-\infty}^{+\infty} \psi^*(x') \delta(-x' - x) \, dx' \right)^* = \\
&= \int_{-\infty}^{+\infty} \psi(x') \delta(-x' - x) \, dx' = \\
&= \psi(-x)
\end{align*}
quindi \( \mathcal{P} \) è hermitiano \( \Rightarrow \mathcal{P} \ket{x} = \ket{-x} \Rightarrow \mathcal{P}^2 \ket{x} = \ket{x} \Rightarrow \mathcal{P}^2 = \mathbbm{1} \Rightarrow \mathcal{P} \mathcal{P} = \mathcal{P}^\dag \mathcal{P} = \mathbbm{1} \Rightarrow \mathcal{P} \) è unitario. \( \mathcal{P}^2 = \mathbbm{1} \Rightarrow \) autovalori \( \pm 1 \).


\bigskip
\noindent
Le autofunzioni sono
\[
\braket{x|\mathcal{P}|\varphi_{-}} = \begin{cases}
- \varphi_{-}(x) \\
\varphi_{-}(-x)
\end{cases} \quad \Rightarrow \quad \varphi_{-} \text{ funzione dispari}
\]
\[
\braket{x|\mathcal{P}|\varphi_{+}} = \begin{cases}
\varphi_{+}(x) \\
\varphi_{+}(-x)
\end{cases} \quad \Rightarrow \quad \varphi_{+} \text{ funzione pari}
\]
\[
f(x) = \frac{f(x) + f(-x)}{2} + \frac{f(x) - f(-x)}{2}.
\]